%%--------------------------------------------------------------------------
%%  SECTION 1: INTRODUCTION AND VISION
%%--------------------------------------------------------------------------
\section{Introduction and Research Vision}

We stand at a remarkable inflection point in computational science.  The most
pressing open problems of our time---controlled nuclear fusion, precision
subsurface imaging for climate-critical carbon sequestration, next-generation
semiconductor design, and predictive modeling of radiative transport in
astrophysical and biomedical settings---share a common mathematical core: they
require solving kinetic equations and associated inverse problems at scales,
speeds, and levels of physical fidelity that are fundamentally inaccessible to
classical numerical methods.

Kinetic theory provides the most complete description of matter at the
mesoscale, the regime bridging quantum mechanics and macroscopic fluid
dynamics.  The Boltzmann equation, the radiative transfer equation (RTE), and
the Vlasov--Maxwell system underpin models of plasma confinement in tokamaks,
photon transport in stellar atmospheres and radiation oncology planning, and
carrier dynamics in semiconductor devices.  These equations define
distributions over phase space of dimension $2d$ ($d = 3$ in three spatial
dimensions), rendering direct discretization intractable for realistic
problems.  Compounding the curse of dimensionality is the ubiquitous presence
of multiple temporal and spatial scales---from the mean free path of a single
particle collision to macroscopic fluid motion---that force classical solvers
to use extremely fine computational grids even when the macroscopic solution
is globally smooth.

Deep learning has demonstrated a striking capacity to approximate
high-dimensional functions and solution operators with sample complexity far
better than classical worst-case theory would predict.  Yet the na\"{i}ve
application of neural networks to physics problems is beset by fundamental
limitations: models trained without respect for physical structure violate
conservation laws, behave catastrophically in stiff multiscale regimes where
the Knudsen number $\varepsilon \ll 1$, and offer no theoretical guarantees
for out-of-distribution generalization.

My research program addresses these limitations through what I call
\emph{full-stack AI for Science}: a vertically integrated methodology that
combines (i)~rigorous mathematical analysis of the underlying physics and the
learning algorithm, (ii)~neural network architectures that encode physical
structure as \emph{hard constraints} rather than soft regularization
penalties, (iii)~scalable implementation exploiting modern accelerated
hardware, and (iv)~production-quality open-source software that makes these
methods accessible to the broader scientific community.  This program has
produced more than 30 peer-reviewed articles and a suite of open-source
packages---hosted at
\href{https://github.com/mazhengcn}{\texttt{github.com/mazhengcn}}---that
are actively used by research groups in applied mathematics, physics, and
engineering.  The work spans four interconnected thrusts: \textbf{(1)
structure-preserving AI for multiscale kinetic equations}; \textbf{(2) a
mathematical theory of deep neural network training}; \textbf{(3)
AI-driven methods for PDE-constrained inverse problems}; and \textbf{(4)
fast spectral algorithms} that serve as both the theoretical bedrock and the
benchmark gold standard for the AI-based methods.
